    \chapter{KESIMPULAN DAN SARAN}

\section{Kesimpulan}
Berdasarkan hasil data pengujian dan analisis terhadap hasil sesungguhnya, dapat disimpulkan bahwa kit eksperimen Pergeseran Wien berbasis digitalisasi Arduino Uno dan sensor TCS34725 \textit{interface} Python 3.8 dapat bekerja dengan optimal dan akurat dalam mengamati hubungan secara eksperimental suhu terhadap panjang gelombang dari radiasi benda hitam berdasarkan Hukum Pergeseran Wien yaitu berbanding terbalik dan berkaitan dengan konstanta pergeseran Wien $(b)$. Berdasarkan hasil pengujian, diperoleh nilai konstanta b sebesar 0,002897773 dengan nilai ketelitian mencapai 1,8796E-09 dan ketepatan mencapai 99,99\% terhadap nilai konstanta $b$ berdasarkan literatur sebesar 0,002898. 

\section{Saran}
Berdasarkan pengujian dan analisis yang telah dilaksanakan, terdapat beberapa saran terkait pengembangan kit eksperimen sebagai berikut.
\begin{enumerate}
    \item Pengembangan alat menggunakan sensor spektrometer yang dapat membaca seluruh spektrum panjang gelombang cahaya tampak hingga inframerah.
    \item Pengembangan visualisasi data dan \textit{interface} dalam tingkat akurasi dan ketelitian pengukuran.
\end{enumerate}
Diharapkan saran tersebut dapat menjadi acuan pengembangan penelitian untuk memperoleh hasil eksperimen yang lebih baik dan optimal.

